\section{Research Methods}
\label{chp:research-methods}

For this proposed research, we envision a multi-method approach that is a combination of these two research methods: (i) Ethnographic Case-Studies (Participant and Non-Participant Observation) and (ii) Multivocal Literature Review.

\subsection{Ethnographic Case-Studies}

An ethnographic case study is a qualitative research approach that combines an in-depth, context-rich investigation of a specific case with ethnographic methods such as participant and non-participant observation, interviews, and document analysis. It seeks to understand social practices, cultural meanings, and lived experiences within a particular setting, emphasizing the relations between individual agency and broader structures~\cite{edmonds2017-research-design}.

As mentioned, the researchers have a vast background in academia and practice with the Linux ecosystem, firsthand or inherited from the experience of other researchers and people involved in the ecosystem. These two sources of insights are what we call \textit{Participant Observation} and \textit{Non-Participant Observation}, respectively, and together represent the primary research method of this project.

\subsubsection{Participant Observation}

In regards to the former, the aim of the participant observation is for the doctoral student to continue his involvement in the Linux project in subsystems like the AMD GPU but go deeper in terms of other subprojects (\textit{Rust for Linux}\footnote{\href{https://rust-for-linux.com/}{Rust for Linux project documentation.}} is one marked to be explored) and in terms of the impact of contributions and possibly escalating to the role of file/driver maintainer. The plan is to systematically catalog the experiences and interactions to build a contributing journal. Also, as noted, the doctoral student maintains two supporting projects for the Linux ecosystem, kw and patch-hub, which are rich sources of insights.

\subsubsection{Non-Participant Observation}

Concerning the latter, the non-participant observation comprises three complementary fronts:

\begin{enumerate}
    \item Mentor new contributors (newcomers) to the Linux ecosystem and collect data about their experience through cataloging mentors' perspectives and observations;
    \item Conduct experiments with new contributors (newcomers) to the Linux ecosystem and, in a separate trial of experiments, veteran kernel developers to qualitatively and quantitatively measure how the FLOSS tools produced cover workflows and mitigate bottlenecks. For the newcomers' trials, questionnaires will be used to collect data, while for the veteran trials, data will be collected through telemetry using an anonymization approach similar to the one used by \textit{Debian Popularity Contest}\footnote{\href{https://popcon.debian.org/}{Debian Popularity Contest website}.};
    \item Interact with key players from the Linux ecosystem and collect feedback through \textit{in-loco} participation in community events, which was a suggestion given by the board of the Linux Foundation to our research group to the detriment of conducting surveys.
\end{enumerate}

\subsection{Multivocal Literature Review}

A \textit{Multivocal Literature Review} (MLR) is a systematic approach to synthesizing knowledge from diverse sources, including academic publications, Grey Literature (GL), practitioner reports, and other non-traditional texts. Essentially combining \textit{Systematic Literature Review} (SLR) and \textit{Grey Literature Review} (GLR), it acknowledges and integrates multiple perspectives and understandings of a research topic while identifying convergence, divergence, and gaps in the discourse~\cite{vahid2018-mlr}. \textit{Wen (2021)} \cite{wen2021-masterthesis}, in her master's thesis, conducted an MLR to identify the gaps between state-of-the-art (academia understanding) and state-of-the-practice (practitioners understanding). As a secondary result, \textit{Wen (2021)} \cite{wen2021-masterthesis} used the lessons learned when conducting the GLR to devise ``guidelines to examine FLOSS projects through its community publications'' due to the lack of a consolidated method to do GLR in this specific context. We plan to make an MLR similar to how \textit{Wen (2021)} \cite{wen2021-masterthesis} did, but limiting the scope to development workflows and going back only 5 years since going back further would result in duplication of work.

The MLR will be meaningful for triangulating findings from the ethnographic case studies - our primary method - with state-of-the-art and state-of-the-practice to eliminate biases and strengthen the research.
